\documentclass[11pt]{article}
\usepackage{fontspec}
\usepackage[french]{babel}
\usepackage{amsmath}
\usepackage{amssymb}
\usepackage{amsthm}
\usepackage{mathrsfs}
\usepackage{tikz-cd}
\usepackage{enumitem}

\newcommand{\printedpage}[1]{\marginpar{\footnotesize p.~#1}}
\newcommand{\msp}[1]{\marginpar{\footnotesize [#1]}}
\newcommand{\Filt}{\mathrm{Filt}}
\newcommand{\Fl}{\mathrm{Fl}}
\newcommand{\Ob}{\mathrm{Ob}}
\newcommand{\card}{\operatorname{card}}
\newcommand{\uHom}{\underline{\mathrm{Hom}}}
\newcommand{\uFl}{\underline{\mathrm{Fl}}}
\newcommand{\into}{\lhook\joinrel\longrightarrow}

\setlength{\parindent}{0pt}

\begin{document}

\printedpage{39}
% head: A. Grothendieck : Dérivateurs, Ch. XV, Édition des Universités de Montpellier II et Paris VII
% foot: Version 5 novembre 2010

\msp{page 67}
L'équivalence de $(a)$ et $(a')$, ou de $(b)$ et $(b')$ est plus ou moins tautologique. L'équivalence de $(a)$ et $(b)$ résulte aussitôt du N.B. inclus dans $1^{\circ}$, compte tenu que les cardinaux $\pi$ tels que les objets de $\Filt^{\pi}$ soient $\pi$-accessibles sont cofinaux dans l'ensemble des cardinaux dans $\mathfrak{U}$. (\textit{Cf.} 3.2, page 16).

\textbf{Proposition 5.13.} \textit{Toute réunion finie de parties L-$\pi$-accessibles de} $\Ob\mathcal{M}$ \textit{est L-$\pi$-accessible. L'intersection d'une famille $(M_{\alpha})_{\alpha\in A}$ de parties L-$\pi$-accessibles, avec} $\card A \leq \pi$\textit{, est L-$\pi$-accessible.}

\textsc{Démonstration}. La première assertion provient du fait qu'une réunion \textit{finie} de parties L-stables de $M$ est L-stable, et qu'on a itou pour la notion de $\pi$-accessibilité. Les deux choses se prouvent d'ailleurs pratiquement de la même façon, en utilisant le

\textbf{Lemme 5.14.} \textit{Soit $I$ un ensemble ordonné filtrant, $(I_{\alpha})_{\alpha\in A}$ une famille} finie \textit{non vide de parties de $I$, telle que $\bigcup I_{\alpha} = I$. Alors il existe $\alpha \in A$ tel que $I_{\alpha}$ soit cofinal.}

\msp{page 68}
\textsc{Démonstration} par l'absurde. Si pour aucun $\alpha$, $I_{\alpha}$ n'était cofinal, on en conclurait que pour tout $\alpha$ il existe un $i_{\alpha}$ tel que $I/i_{\alpha} \cap I_{\alpha} = \emptyset$. Soit $i$ majorant les $i_{\alpha}$, on aura donc $I/i \cap I_{\alpha} = \emptyset \ \ \forall\, \alpha \in A$, contrairement à l'hypothèse que $\bigcup I_{\alpha} = I$.

Le cas d'une \textit{intersection} de parties L-$\pi$-accessibles : elle est bien sûr L-stable, donc il reste à prouver qu'elle est $\pi$-accessible, \textit{i.e.} satisfait la condition $\mathrm{b}_{\pi}$ de \textit{loc. cit.}

\msp{page 69}
Soit donc
\[
X \;=\; \varinjlim_{i\in I} X_{i} \; , \qquad\quad X_{i} \in \Filt^{\pi} \; ,
\]
une représentation $\pi$-adaptée de $X$, avec $I$ un ensemble ordonné $\pi$-adapté. Donc on a
\[
\begin{aligned}
X \in M_{\alpha} \quad &\Longleftrightarrow \quad \text{l'ensemble } I_{M_{\alpha}} \text{ des } i \in I \text{ tels que } X_{i} \in M_{\alpha}\\
&\phantom{\Longleftrightarrow} \quad \text{est cofinal dans } I.
\end{aligned}
\tag*{$(*)$}
\]
D'ailleurs, comme $M_{\alpha}^{\pi} = M_{\alpha} \cap \Filt^{\pi}$ est stable par $\varinjlim$ filtrante de cardinal $\leq \pi$, on voit tout de suite que la partie $I_{M_{\alpha}}$ de $I$ est stable par bornes supérieures dans $I$ de parties $J$ de $I_{M_{\alpha}}$ filtrantes et de cardinal $\leq \pi$. Donc $I_{M_{\alpha}}$ cofinal équivaut à $I_{M_{\alpha}}$ $\pi$-\textit{cofinal} (\textit{cf.} 5.8, page 59). Soit $M = \bigcap_{\alpha\in A} M_{\alpha}$, je veux prouver que
\[
X \in M \quad \Longleftrightarrow \quad \text{l'ensemble } I_{M} \ (= \textstyle\bigcap I_{M_{\alpha}}) \text{ est cofinal dans } I.
\tag*{$(**)$}
\]
Or $(*)$ implique que $X \in M$ $(= \bigcap M_{\alpha})$ équivaut à $I_{M_{\alpha}}$ $\pi$-cofinal pour tout $\alpha$. Donc au total on veut prouver
\[
I_{M_{\alpha}} \text{ cofinal dans } I \text{ pour tout } \alpha \quad\ \Longleftrightarrow \quad\ \textstyle\bigcap I_{M_{\alpha}} \text{ est cofinal dans } I.
\tag*{$(***)$}
\]

\msp{page 70}
On a bien sûr $\Longleftarrow$, donc il reste à prouver $\Longrightarrow$. Si l'ensemble d'indices $A$ est fini, cela résulte du lemme 5.8, compte tenu que la cofinalité des $I_{M_{\alpha}}$ implique leur $\pi$-cofinalité. On doit donc simplement améliorer 5.8, par le

\printedpage{40}
% head: A. Grothendieck : Dérivateurs, Ch. XV, Édition des Universités de Montpellier II et Paris VII
% foot: Version 5 novembre 2010

\textbf{Lemme 5.15.} \textit{Soient $\pi$ un cardinal, $I$ un ensemble ordonné filtrant $\pi$-adapté, et soit $(I_{\alpha})_{\alpha\in A}$ une famille de parties $I_{\alpha}$ $\pi$-cofinales} (cf. 5.8), \textit{telle que} $\card A \leq \pi$. \textit{Alors $\bigcap I_{\alpha}$ est $\pi$-cofinale.}

Il suffit de prouver qu'elle est cofinale, car la propriété de stabilité par bornes supérieures de parties filtrantes de cardinal $\leq \pi$ est évidente. Quitte à remplacer $I$ par $I/i_{0}$, et les $I_{\alpha}$ par leurs traces sur $I/i_{0}$, la question est simplement de prouver que $\bigcap_{\alpha} I_{\alpha} \neq \emptyset$. On va construire un élément $s$ de $\bigcap I_{\alpha}$ par induction transfinie. Prenons un bon ordre sur $A$, et prouvons par récurrence transfinie
\msp{page 71}
que pour tout $\alpha \in A$, l'ensemble
\[
J_{\alpha} \;=\; \bigcap_{\beta<\alpha} I_{\beta}
\]
est cofinal. Si on sait que $J_{\alpha}$ est cofinal, il en est de même de $J_{\alpha+1} = J_{\alpha} \cap I_{\alpha}$, par 5.8. Donc on est ramené à prouver que $J_{\alpha}$ est cofinal si $\alpha$ est un ordinal limite, donc à prouver le lemme 5.15 pour une suite transfinie \textit{décroissante} de parties $\pi$-cofinales $J_{\beta}$ $(\beta < \alpha)$ quand $\card \alpha \leq \pi$, telle que $J_{\beta} = \bigcap_{\beta'<\beta} J_{\beta'}$ quand $\beta$ est limite. À prouver que l'intersection est $\pi$-cofinale -- et on est réduit à prouver qu'elle est non vide. On va construire par récurrence transfinie une suite \textit{croissante} d'éléments
\[
i_{\beta} \;\in\; J_{\beta} \qquad\quad (\, \beta \;\leq\; \alpha \,) \; .
\]
Si $i_{\beta}$ est construit, comme $J_{\beta+1}$ est cofinal, il existe $i_{\beta+1} \geq i_{\beta}$ dans $J_{\beta+1}$, O.K. Si $\beta$ est un ordinal limite, on prend
\[
i_{\beta} \;=\; \sup_{\beta'<\beta} i_{\beta'} \; ,
\]
$i_{\beta}$ est dans chaque $I_{\beta_{0}}$ pour $\beta_{0} < \beta$,
\msp{page 72}
puisque les $i_{\beta'}$ pour $\beta_{0} \leq \beta' < \beta$ sont tous dans $J_{\beta_{0}} \supset J_{\beta'}$, donc aussi leur limite, donc
\[
i_{\beta} \;\in\; \bigcap_{\beta_{0}<\beta} J_{\beta_{0}} \;=\; J_{\beta} \; .
\]
Donc la récurrence transfinie marche, et nous donne donc un élément $s_{\alpha}$ de $J_{\alpha}$\,!

Je vais prouver aussi en forme l'assertion pour $\bigcup_{\alpha\in A} M_{\alpha}$, quand $A$ est fini. Posant cette fois $M = \bigcup_{\alpha} M_{\alpha}$, il faut prouver, moyennant les relations $(*)$ (p. 68), qu'on a
\[
\begin{aligned}
&X \in M &&\Longleftrightarrow \quad &&\text{l'ensemble } \ I_{M} \ \ (= \ \textstyle\bigcup_{\alpha} I_{M_{\alpha}})\\
&(\textit{i.e.} \ \exists \ \alpha \text{ avec } X \in M_{\alpha}) && &&\text{est cofinal dans } I \ ,
\end{aligned}
\tag*{$(**)$}
\]
et moyennant $(*)$ le premier terme de cette équivalence à établir signifie qu'il existe $\alpha$ avec $I_{M_{\alpha}}$ cofinal dans $I$. Bien sûr, si un des $I_{M_{\alpha}}$ est cofinal dans $I$, a fortiori leur réunion. Mais inversement, si la réunion est cofinale, elle est filtrante et réunion des parties (filtrantes) $I_{M_{\alpha}}$, dont l'une doit donc être cofinale dans elle par le lemme 5.14, donc aussi dans $I$, \textit{q.e.d.}

\msp{page 73}
\textbf{Corollaire 5.16} (de la proposition 5.13). \textit{Toute réunion finie de parties L-accessibles de $\mathcal{M}$ est L-accessible. Toute intersection $\bigcap M_{\alpha}$ d'une famille petite de parties L-accessibles de $\mathcal{M}$, est L-accessible.}

\printedpage{41}
% head: A. Grothendieck : Dérivateurs, Ch. XV, Édition des Universités de Montpellier II et Paris VII
% foot: Version 5 novembre 2010

\section*{6\quad Essai avorté pour un critère d'engendrement des parties L-accessibles de $\Fl(\mathcal{M})$.}

Soit $\mathcal{M}$ une catégorie munie d'une filtration cardinale $(\Filt^{\pi})_{\pi\geq\pi_{0}}$, et soit $S$ un objet de Cat, considérons
\[
\mathcal{M}(S) \;=\; \uHom(S,\mathcal{M}) \; ,
\]
et munissons $\mathcal{M}(S)$ de la filtration par les
\[
\mathcal{M}^{\pi}(S) \;=\; \uHom(S,\mathcal{M}^{\pi} = \underline{\Filt}^{\pi}(\mathcal{M})) \; .
\]
Je dis que c'est là une filtration cardinale de $\mathcal{M}(S)$. Il faut vérifier les conditions $(a)$, $(b)$, $(c)$ de SGA 4 I 9.12. Les conditions $(a)$ ($\Filt^{\pi}$ essentiellement petite) et $(b)$ (condition $\mathrm{L}_{\pi_{0}}$, et stabilité de $\Filt^{\pi}$ par $\varinjlim\limits_{I}$ , $I$ filtrante grande devant $\pi_{0}$ et de cardinal $\leq \pi$) passent aussitôt de $\mathcal{M}$ à $\mathcal{M}(S)$. La condition $(c)$ est plus délicate (représentation
\msp{page 74}
d'un objet comme $\varinjlim$ filtrante grande devant $\pi$ d'objets de $\Filt^{\pi}$). Je ne vais pas refaire le travail, fait dans un cas beaucoup plus général et plus délicat dans SGA 4 I §9, savoir dans le corollaire 9.24 du théorème 9.22 (p. 177). La seule restriction, c'est qu'on y suppose les objets de $\mathcal{M}$ accessibles, et désignant par $\pi_{1}$ le plus petit cardinal tel que $\pi_{1} \geq \pi_{0}$ et que les objets de $\Filt^{\pi_{0}}$ soient $\pi_{1}$-accessibles, on se borne pour la filtration cardinale de $\uHom(S,\mathcal{M})$ aux cardinaux \fbox{$\pi \geq c = 2^{c_{0}}$,} où
\[
c_{0} \;=\; \sup(\pi_{1}, \card \Fl(S)) \; .
\]
C'est une restriction raisonnable et anodine.

Pour l'instant, je suis intéressé surtout au cas où $S = \boldsymbol{\Delta}_{1}$, donc à
\[
\uHom(\boldsymbol{\Delta}_{1},\mathcal{M}) \;=\; \uFl(\mathcal{M}) \; ,
\]
auquel cas on a $c_{0} = \pi_{1}$, donc il suffit de se borner aux cardinaux $\pi \geq 2^{\pi_{1}}$ pour avoir une bonne filtration cardinale sur $\uFl(\mathcal{M})$.

C'est donc l'ensemble des deux conditions $(b)$, $(c)$ sur $\mathcal{M}$ (p. 44) qui est stable par passage aux $\uHom(S,\mathcal{M}) = \mathcal{M}(S)$. Il en est de même de la condition $(a)$ (stabilité par petites limites inductives), et aussi de la condition\msp{page 75}
$(d)$ (p. 45). De même la condition d'exactitude $(e)$ (p. 46) sur les $\varinjlim$ filtrantes (elles doivent transformer les monomorphismes en monomorphismes).

Pour être tranquille, je vais faire sur $\mathcal{M}$ les hypothèses $(a)$ à $(e)$ inclus (bien que $(e)$ me semble un peu artificielle, on devrait pouvoir l'éliminer, dans un second souffle\,!). Ces conditions ont donc la propriété de stabilité souhaitable par passage aux $\mathcal{M}(S)$, et en particulier par passage à $\Fl(\mathcal{M})$. Quitte à augmenter $\pi_{0}$, on peut donc supposer que les $\Fl^{\pi}(\mathcal{M}) = \mathrm{Hom}(\boldsymbol{\Delta}_{1},\mathcal{M}^{\pi})$ $(\pi \geq \pi_{0})$ forment aussi une filtration cardinale de $\uFl(\mathcal{M})$.

On dispose donc des notions de parties $\pi$-accessibles, accessibles, L-$\pi$-accessibles, L-accessibles de $\Fl(\mathcal{M}) \ = \ \Ob\uFl(\mathcal{M})$. Ces notions sont stables par réunions et intersections finies, et les notions absolues (sans $\pi$) stables par intersections ``quelconques'' (\textit{i.e.} pour des \textit{petites} familles de parties). Ce sera là une chose particulièrement\printedpage{42}%
% head: A. Grothendieck : Dérivateurs, Ch. XV, Édition des Universités de Montpellier II et Paris VII
% foot: Version 5 novembre 2010
\msp{page 76}
cruciale, le cas d'une intersection $C \cap W$\,!

Le résultat principal, pour lequel j'ai fait les gammes du paragraphe précédent, est le suivant :

\setcounter{footnote}{25}
\textbf{Théorème 6.1}\footnote{Démonstration canulée, énoncé peut-être \textit{faux}\,!!}. \textit{Soit $\mathcal{M}$ une catégorie satisfaisant les conditions dites :}
\begin{enumerate}[label={}]
\item[a)] \textit{Stabilité par petites} $\varinjlim$ \textit{(pas nécessairement filtrantes).}
\item[b)] $\mathcal{M}$ \textit{munie d'une filtration cardinale} $(\Filt^{\pi})_{\pi\geq\pi_{0}}$.
\item[c)] \textit{Les objets de $\mathcal{M}$ sont accessibles.}
\item[d)] \textit{Les} $\Filt^{\pi}$ \textit{sont stables par} $\varinjlim$ \textit{indexées par des catégories $I$ telles que} $\card \Fl(I) \leq \pi$\textit{, et si $Z \in \Filt^{\pi}$, tout quotient $Z'$ de $Z$ (tel que $Z \longrightarrow Z'$ ne se factorise pas par un sous-objet $Z'_{0}$ de $Z'$, différent de $Z'$) est dans} $\Filt^{\pi}$.
\item[e)] \textit{Les limites inductives filtrantes dans $\mathcal{M}$ transforment les monomorphismes en monomorphismes.}
\end{enumerate}
\textit{Soit $C$ une partie L-accessible de} $\Fl(\mathcal{M}) = \Ob\uFl(\mathcal{M})$, i.e. (5.12) \textit{une partie stable par} $\varinjlim$ \textit{filtrantes, et accessible. Supposons de plus $C$ formée de monomorphismes}\footnote{Finalement, on se débarrasse de cette hypothèse.}\textit{, et que $C$ contient les} isomorphismes. \textit{Alors il existe une partie petite $C_{0} \subset C$, telle que $C \subset \widetilde{C_{0}}$. Plus précisément, si $\pi$ est un cardinal tel que $C$ soit $\pi$-accessible, et si on pose $C^{\pi} = C \cap \Fl^{\pi}(\mathcal{M})$, alors on a}
\msp{page 77}
\[
C \;\subset\; \widetilde{C^{\pi}} \;=\; \widetilde{C_{0}} \; ,
\tag*{(6.1.1)}
\]
\textit{où $C_{0}$ est n'importe quelle petite partie de $C^{\pi}$ telle que tout objet de $C^{\pi}$ soit isomorphe à un objet de $C_{0}$.}

\textbf{Corollaire 6.2.} \textit{On a $C_{*} = C_{0\,*}$, $\widetilde{C} = \widetilde{C_{0}}$.}

\textbf{Corollaire 6.3.} \textit{Conditions équivalentes sur la partie L-accessible $C$ de} $\Fl(\mathcal{M})_{\mathrm{mono}}$\textit{, contenant les isomorphismes :}
\begin{enumerate}[label={}]
\item[a)] \textit{$C = \widetilde{C}$.}
\item[b)] \textit{$C$ est stable par cochangement de base et par composition.}
\item[b$'$)] \textit{$C$ est stable par cochangement de base, et par sommes amalgamées simples,} i.e. \textit{si $A \longrightarrow B_{1}$, $A \longrightarrow B_{2}$ sont dans $C$, alors $A \longrightarrow B_{1} \amalg_{A} B_{2}$ aussi.}
\item[c)] \textit{$C^{\pi}$ est stable dans} $\Fl^{\pi}(\mathcal{M}) = \Fl(\mathcal{M}^{\pi})$ \textit{par cochangement de base et par composition.}
\item[c$'$)] \textit{$C^{\pi}$ est stable dans} $\Fl(\mathcal{M}^{\pi})$ \textit{par cochangement de base et par sommes amalgamées} (cf. b$'$).
\end{enumerate}
\textit{Quand ces conditions sont satisfaites, alors $(C,C_{*})$ est un couple de Quillen, justiciable du théorème} 1 \textit{du paragraphe} 1 (p. 10).

La dernière assertion résulte du théorème 6.1 et du corollaire au théorème 1 du paragraphe 1. On a
\printedpage{43}%
% head: A. Grothendieck : Dérivateurs, Ch. XV, Édition des Universités de Montpellier II et Paris VII
% foot: Version 5 novembre 2010
\msp{page 78}
d'autre part tautologiquement
\[
\begin{tikzcd}
(a) \arrow[r, Rightarrow] & (b) \arrow[r, Rightarrow] \arrow[d, Rightarrow] & (b') \arrow[d, Rightarrow] \\
 & (c) \arrow[r, Rightarrow] & (c') \ ,
\end{tikzcd}
\]
et il reste à prouver $(c') \Longrightarrow (a)$, ce qu'on fera plus bas.

Prouvons donc 6.1. Dans (6.1.1) l'égalité $\widetilde{C_{0}} = \widetilde{C^{\pi}}$ est tautologique. Donc tout revient à prouver
\[
C \;\subset\; \widetilde{C^{\pi}} \; .
\]
Une difficulté ici, c'est qu'on ne sait \textit{pas} que $\widetilde{C^{\pi}}$ soit stable par $\varinjlim$ filtrantes -- on le sait seulement pour $C$. Soit donc
\[
A \;\into\; B
\]
une flèche dans $C$, prouvons qu'elle est dans $\widetilde{C^{\pi}}$. Pour ceci, on va construire, par induction transfinie, une suite transfinie \textit{strictement} croissante de sous-objets $A_{\alpha}$ de $B$ contenant $A = A_{0}$,
\[
\underbrace{A}_{=\,A_{0}} \;\; \overset{i_{\alpha}}{\into} \;\; A_{\alpha} \;\; \overset{i_{\beta,\alpha}}{\into} \;\; A_{\beta} \xrightarrow{\ j_{\beta}\ } B \; ,
\]
telle que l'on ait
\[
i_{\alpha} \;\in\; \widetilde{C^{\pi}} \; , \qquad\quad \text{plus généralement les } i_{\beta,\alpha} \;\in\; \widetilde{C^{\pi}} \; , \qquad\quad j_{\alpha} \;\in\; C \; ,
\]
et que l'on ait
\[
A_{\alpha} \;=\; \varinjlim_{\beta<\alpha} A_{\beta} \qquad\quad \text{si } \alpha \text{ ordinal limite} \; .
\]

\msp{page 79}
Cela précise déjà la définition récurrente pour $\alpha = \alpha_{0}$ ordinal limite. Notons que le $A_{\alpha_{0}}$ construit en termes des $A_{\beta}$ $(\beta < \alpha_{0})$ satisfait bien aux conditions que
\begin{enumerate}[label={}]
\item[\textit{a})] les $i_{\alpha_{0},\beta} : A_{\beta} \longrightarrow A_{\alpha_{0}}$ sont dans $\widetilde{C^{\pi}}$, car $\widetilde{C^{\pi}}$ est stable par ``limites inductives ordinales strictes'' (\textit{cf.} 1.19, p. 8),
\item[\textit{b})] $j_{\alpha_{0}} \, : \, A_{\alpha_{0}} \longrightarrow B$ est dans $C$, car c'est la limite inductive filtrante du système des $A_{\beta} \longrightarrow B$ $(\beta < \alpha_{0})$ qui sont dans $C$, et $C$ est stable par ces limites.
\end{enumerate}
Il reste donc à préciser comment on passe de $A_{\alpha}$ à $A_{\alpha+1}$. Cela revient à prouver le

\textbf{Lemme 6.4.} \textit{Soit $f : A \into B$ dans $C$, telle que $f$ ne soit pas un isomorphisme. Alors il existe un sous-objet $A_{1}$ de $B$,} distinct de $A$, \textit{tel que $A \longrightarrow A_{1}$ soit dans $\widetilde{C^{\pi}}$, et $A_{1} \longrightarrow B$ dans $C$.}

Si ce lemme est admis, cela signifie que la construction transfinie se continue aussi longtemps que l'on trouve des $A_{\alpha}$ tels que $A_{\alpha} \neq B$, \textit{i.e.} $A_{\alpha} \into B$ ne soit pas un isomorphisme. Mais comme les $A_{\alpha}$ sont tous différents, si $\pi$ est l'ensemble des sous-objets de $B$ (on sait que c'est un

\end{document}
